%% COVER LETTER, uploaded to Editorial Manager as a separate file.
%% Fill in the bracketed parts. Keep it to one page.
\documentclass[11pt,a4paper]{article}
\usepackage[margin=2.5cm]{geometry}
\usepackage[T1]{fontenc}
\usepackage[utf8]{inputenc}
\usepackage{mathptmx}
\usepackage[hidelinks]{hyperref}
% The one file to edit before submitting: every edition reads it.
\newcommand{\capmauthor}{Henrique Hoffmann Stefanutti}
\newcommand{\capmaffiliation}{Independent Researcher}
\newcommand{\capmcity}{S\~ao Carlos}
\newcommand{\capmcountry}{Brazil}
\newcommand{\capmemail}{henrique407@outlook.com}
\newcommand{\capmorcid}{0000-0002-1630-1655}
\newcommand{\capmrepo}{https://github.com/rikosoo/InfraCriticalCar}

\pagestyle{empty}

\begin{document}

\noindent \capmauthor\\
\capmaffiliation\\
\capmcity, \capmcountry\\
\capmemail

\vspace{0.6cm}
\noindent \today

\vspace{0.6cm}
\noindent To the Editor-in-Chief\\
\textit{International Journal of Critical Infrastructure Protection}

\vspace{0.6cm}

\noindent Dear Editor,

\vspace{0.3cm}

I am submitting the manuscript ``From Cloud to Assembly Line: Modelling
Cyberattack Paths in Automotive Manufacturing as Critical Infrastructure'' for
consideration as a research paper in the \textit{International Journal of
Critical Infrastructure Protection}.

\vspace{0.3cm}

\noindent\textbf{Fit with the journal's scope.} The paper treats the automotive
assembly plant, rather than the vehicle, as critical infrastructure. It models
how a compromise propagates across corporate information technology, identity,
cloud services, the industrial demilitarised zone and the control layers, and
how a failure at one supplier cascades to an original equipment manufacturer
through just-in-time coupling --- the interdependency and cascading-failure
problem the journal addresses, in a sector that regulators have recently brought
within critical-infrastructure obligations. The work combines an engineering
model with the regulatory language plants actually use, mapping every control to
IEC 62443 system requirements and to the NIST Cybersecurity Framework.

\vspace{0.3cm}

\noindent\textbf{Contribution.} The paper contributes (i) a curated corpus of 20
publicly reported incidents in automotive manufacturing from 2017 to 2025, each
reconstructed as an asset-level path with an explicit confidence label; (ii) a
cross-plane attack-path model with an explicit probability semantics, instantiated
as a reference plant of 12 zones, 48 assets and 91 attack steps; (iii) a
quantitative analysis combining exact most-likely-path ranking with Monte Carlo
campaign simulation; and (iv) a comparison against four baseline prioritisations,
including a result against our own model that we report rather than omit. The
model, corpus and experiment pipeline are released openly, and every number in
the paper is regenerated by a single reproducible run.

\vspace{0.3cm}

\noindent\textbf{Originality and prior submission.} The manuscript is original,
has not been published elsewhere and is not under consideration by any other
journal. % If it was previously submitted and rejected elsewhere, say so here,
% summarise the changes made, and attach the previous decision letter.

\vspace{0.3cm}

\noindent\textbf{Generative AI.} [State here any use of generative artificial
intelligence in preparing the manuscript, the tool and the purpose, consistent
with the declaration in the manuscript.]

\vspace{0.3cm}

\noindent\textbf{Suggested reviewers.} [Optional: three to five names with
affiliations and e-mail addresses, none of them recent collaborators.]

\vspace{0.6cm}

\noindent Thank you for considering this submission.

\vspace{0.6cm}

\noindent Yours sincerely,\\
\capmauthor

\end{document}
